\subsection{Additional evaluation criteria}\label{app:A}
\subsubsection*{Equations}\label{app:A:equations}
In the following equations $y_t$ and $\hat{y}_t$ is the ground truth target 
value at prediction time and the predicted value at prediction time 
respectively. \\
\Cref{eq:rmse} depicts how to calculate the \gls{rmse}.
\begin{equation}
	\text{RMSE} = \sqrt{\frac{1}{n} \sum_{t=1}^{n} (y_t - \hat{y}_t)^2}
	\label{eq:rmse}
\end{equation}
This metric is similar to \gls{mae}, however, it squares the errors first, to 
punish outliers more severely. In a market with frequent spikes, both positive 
and negative, this behaviour may not be ideal. While it could have been used as 
a performance metric, \gls{mae} was chosen instead due to their inherent 
similarities and its easier interpretability. \\

\Cref{eq:r2} depicts how to calculate the \gls{r2}.
\begin{equation}
	R^2 = 1 - \frac{\sum_{t=1}^{n} (y_t - \hat{y}_t)^2}{\sum_{t=1}^{n} (y_t - 
		\bar{y})^2}
	\label{eq:r2}
\end{equation}
This metric can be viewed in the same sense as an accuracy score in binary 
classification. As is the case with accuracy in such scenarios, it gives you a 
fair wholistic view of the models, but may be ineffective in capturing the 
intricacies of what you are actually aiming for, and therefore this metric was 
not chosen. \\

\Cref{eq:wmape} depicts how to calculate the \gls{wmape}.
\begin{equation}
	\text{WMAPE} = \frac{\sum_{t=1}^{n} |y_t - \hat{y}_t|}{\sum_{t=1}^{n} |y_t|}
	\label{eq:wmape}
\end{equation}
This metric represents the errors as a percentage, which is easily 
understandable, however, in the case of energy markets, where prices may hover 
around zero, or in the negative values, for extended periods of time, this 
metric may become unstable, due to the low numbers involved. \\

\Cref{eq:smape} depicts how to calculate the \gls{smape}. 
\begin{equation}
	\text{sMAPE} = \frac{1}{n} \sum_{t=1}^{n} \frac{2|y_t - \hat{y}_t|}{|y_t| + 
		|\hat{y}_t|}
	\label{eq:smape}
\end{equation}
This metric is similar to \gls{wmape}, however it treats positive and negative 
errors as equally problematic, which is good in the case of price prediction, 
since both high-balling and low-balling the price of energy is problematic when 
trading energy. However, it has the some problem as \gls{wmape}, of becoming 
unstable during low-number periods. \\

\Cref{eq:mda} depicts how to calculate the \gls{mda}.
\begin{equation}
	\text{MDA} = \frac{1}{n-1} \sum_{t=1}^{n-1} \mathbf{1}_{\text{sgn}(y_{t+1} 
		- y_t) == \text{sgn}(\hat{y}_{t+1} - y_t)}
	\label{eq:mda}
\end{equation}
This metric signifies what percentage of the time the model was correct in the 
direction of the next price, meaning, did the price go up, or did the price go 
down. This is useful for energy traders, however, directionality itself is not 
of a measurable magnitude, and does not align with the premise of this 
comparison of how data constellations and horizons impact model behaviour, 
therefore, this metric was not chosen for this paper. 